\section{Modelos Teóricos}

\begin{figure}[!h]
  \centering
  \begin{tikzpicture}
    \coordinate (A) at (0,0);
    \coordinate (C) at (3,3);
    \coordinate (B) at (5,-1);
    \coordinate (h) at (1,0);
    \coordinate (nh) at ($(B)+(-1,0)$);
    \coordinate (v) at (0,1);
    \fill[color=black] (A) circle [radius=3pt] node[left] {$A$};
    \fill[color=black] (C) circle [radius=3pt] node[above] {$C$};
    \fill[color=black] (B) circle [radius=3pt] node[right] {$B$};
    \draw[ultra thick] (A) -- (C) node[midway, above left] {$a$};
    \draw[ultra thick] (B) -- (C) node[midway, above right] {$a$};
    \pic [draw, angle radius=0.8cm] {angle = h--A--C};
    \node at (A) [shift={(25:1cm)}] {$\beta$};
    \draw (A) -- (h);
    \pic [draw, angle radius=0.8cm] {angle = C--B--nh};
    \node at (B) [shift={(145:1cm)}] {$\gamma$};
    \draw (B) -- (nh);
    \draw[thick, ->] (A) -- +(0,1) node[above] {$V_A$};
    \draw[thick, ->] (B) -- +(0,-1) node[below] {$V_B$};
    \draw[thick, ->] (C) -- +(1,0.5) node[above] {$V_C$};
  \end{tikzpicture}
  \caption{Diagrama cinemático del problema.}
\end{figure}

El análisis cinemático del mecanismo se lleva a cabo mediante el método de la
velocidad relativa, el cual permite establecer relaciones vectoriales entre los
distintos eslabones del sistema \cite{hibbeler2016}. Este método resulta
especialmente adecuado para mecanismos de cuatro barras, ya que descompone el
movimiento absoluto de un punto en términos del movimiento de un punto de
referencia más el movimiento relativo respecto a dicho punto
\cite{meriam2018}. Para resolver las velocidades se plantean dos ecuaciones
vectoriales con dos incógnitas, una por cada brazo del mecanismo. Con este
propósito, se construyen los diagramas de velocidad relativa para cada brazo,
comenzando por el brazo $AC$.

\begin{figure}[!h]
  \centering
  \begin{tikzpicture}
    \coordinate (A) at (0,0);
    \coordinate (C) at (3,3);
    \coordinate (h) at (1,0);
    \coordinate (v) at (0,1);
    \fill[color=black] (A) circle [radius=3pt] node[left] {$A$};
    \fill[color=black] (C) circle [radius=3pt] node[above right] {$C$};
    \draw[ultra thick] (A) -- (C) node[midway, above left] {$a$};
    \pic [draw, angle radius=0.8cm] {angle = h--A--C};
    \node at (A) [shift={(25:1cm)}] {$\beta$};
    \draw (A) -- (h);
    \draw[thick, ->] (A) -- +(0,1) node[above] {$V_A$};
    \draw[thick, ->] (C) -- +(1,-1) node[right] {$\vec{V}_{C/A}$};
    \draw ($(C)+(-0.3,-0.3)$) -- +(0.3,-0.3) -- +(0.6,0);
    \coordinate (Cv) at ($(C)+(0,-1)$);
    \coordinate (Cv2) at ($(C)+(1,-1)$);
    \pic [draw, angle radius=0.8cm] {angle = Cv--C--Cv2};
    \node at (C) [shift={(-70:1.1cm)}] {$\beta$};
    \draw (C) -- +($(0,-1)$);
  \end{tikzpicture}
  \caption{Diagrama de velocidad relativa $\vec{V}_{C/A}$ respecto al brazo $AC$.}
\end{figure}

Dado que la velocidad relativa $\vec{V}_{C/A}$ es perpendicular al brazo $AC$
\cite{beer2019}, su dirección forma un ángulo $\beta$ con la vertical, tal como
se ilustra en la figura anterior. Análogamente, se construye el diagrama
correspondiente al brazo $CB$.

\begin{figure}[!h]
  \centering
  \begin{tikzpicture}
    \coordinate (B) at (5,-1);
    \coordinate (C) at (3,3);
    \coordinate (nh) at ($(B)+(-1,0)$);
    \fill[color=black] (B) circle [radius=3pt] node[right] {$B$};
    \fill[color=black] (C) circle [radius=3pt] node[above left] {$C$};
    \draw[ultra thick] (B) -- (C) node[midway, above right] {$a$};
    \pic [draw, angle radius=0.8cm] {angle = C--B--nh};
    \node at (B) [shift={(145:1cm)}] {$\gamma$};
    \draw (B) -- (nh);
    \draw[thick, ->] (B) -- +(0,-1) node[below] {$V_B$};
    \draw[thick, ->] (C) -- +(-1,-0.5) node[left] {$\vec{V}_{C/B}$};
    \draw ($(C)+(0.3,-0.5)$) -- +(-0.3,-0.15) -- +(-0.6,0.3);
    \coordinate (Cv) at ($(C)+(0,-1)$);
    \coordinate (Cv2) at ($(C)+(-1,-0.5)$);
    \pic [draw, angle radius=0.8cm] {angle = Cv2--C--Cv};
    \node at (C) [shift={(250:1.1cm)}] {$\gamma$};
    \draw (C) -- +(0,-1);
  \end{tikzpicture}
  \caption{Diagrama de velocidad relativa $\vec{V}_{C/B}$ respecto al brazo $CB$.}
\end{figure}

Por la misma razón geométrica, $\vec{V}_{C/B}$ es perpendicular al brazo $CB$
y forma un ángulo $\gamma$ con la vertical. A partir de ambos diagramas es
posible expresar las velocidades relativas en notación vectorial cartesiana
\cite{hibbeler2016}:

% ── Velocidad de C respecto a A y B ─────────────────────────────────────────
\begin{equation*}
    \vec{V}_{C/A} = V_{C/A}\sin\beta\,\ihat - V_{C/A}\cos\beta\,\jhat
\end{equation*}

\begin{equation*}
    \vec{V}_{C/B} = V_{C/B}\sin\gamma\,\ihat + V_{C/B}\cos\gamma\,\jhat
\end{equation*}

Aplicando la ecuación vectorial de velocidad relativa \cite{meriam2018}, se
establece la siguiente igualdad:

% ── Ecuación vectorial de velocidad ─────────────────────────────────────────
\begin{equation*}
    \vec{V}_{C/A} + \vec{V}_A = \vec{V}_C = \vec{V}_{C/B} + \vec{V}_B
\end{equation*}

Sustituyendo las expresiones vectoriales y separando la ecuación resultante en
sus componentes $x$ e $y$, se obtiene el siguiente sistema de dos ecuaciones
con dos incógnitas, $V_{C/A}$ y $V_{C/B}$:

% ── Sistema 2×2 para velocidades ────────────────────────────────────────────
\begin{equation*}
    \begin{cases}
        V_{C/A}\sin\beta - V_{C/B}\sin\gamma = 0 \\
        V_{C/A}\cos\beta + V_{C/B}\cos\gamma = V_A + V_B
    \end{cases}
\end{equation*}

Este sistema lineal admite solución directa, de la cual se determinan las
magnitudes de las velocidades relativas y, por consiguiente, las velocidades
angulares $\omega_{C/A}$ y $\omega_{C/B}$ de cada brazo \cite{beer2019}.
Conocidas estas, se procede al análisis de aceleración, para el cual se
requiere la información de posición de cada brazo expresada en forma vectorial
\cite{hibbeler2016}:

% ── Definición de vectores de posición ──────────────────────────────────────
\begin{align*}
    \vec{r}_{C/A} &= a\cos\beta\,\ihat + a\sin\beta\,\jhat\\
    \vec{r}_{C/B} &= -a\cos\gamma\,\ihat + a\sin\gamma\,\jhat
\end{align*}

Adicionalmente, dado que el movimiento es plano, todas las velocidades y
aceleraciones angulares están orientadas en la dirección del eje $z$
\cite{meriam2018}:

% ── Velocidades y aceleraciones angulares ────────────────────────────────────
\begin{align*}
    \vec{\omega}_{C/A} = \omega_{C/A}\,\khat, \qquad
    \vec{\omega}_{C/B} = \omega_{C/B}\,\khat \\
    \vec{\alpha}_{C/A} = \alpha_{C/A}\,\khat, \qquad
    \vec{\alpha}_{C/B} = \alpha_{C/B}\,\khat
\end{align*}

Con toda esta información, se aplica la ecuación general de aceleración
relativa \cite{beer2019}, la cual incluye tanto el término tangencial,
asociado a la aceleración angular, como el término centrípeto, asociado a la
velocidad angular al cuadrado:

% ── Aceleración de C desde A ─────────────────────────────────────────────────
\begin{align*}
  \vec{a}_C &= \vec{\alpha}_{C/A} \times \vec{r}_{C/A} + \vec{\omega}_{C/A} \times
    \left(\vec{\omega}_{C/A} \times \vec{r}_{C/A}\right) + \vec{a}_A \\
     &= \vec{\alpha}_{C/B} \times \vec{r}_{C/B} + \vec{\omega}_{C/B} \times
    \left(\vec{\omega}_{C/B} \times \vec{r}_{C/B}\right) + \vec{a}_B
\end{align*}

Desarrollando los productos vectoriales e igualando componentes $x$ e $y$ de
ambas expresiones, se obtiene un nuevo sistema lineal de dos ecuaciones con
dos incógnitas, $\alpha_{C/A}$ y $\alpha_{C/B}$, el cual puede escribirse en
forma matricial compacta \cite{hibbeler2016}:

\begin{equation*}
    \resizebox{\columnwidth}{!}{$
    \begin{bmatrix}
        -a\sin\beta & a\sin\gamma \\
        a\cos\beta  & a\cos\gamma
    \end{bmatrix}
    \begin{bmatrix} \alpha_{C/A} \\ \alpha_{C/B} \end{bmatrix}
    =
    \begin{bmatrix}
        a(\omega_{C/B}^{2}\cos\gamma - \omega_{C/A}^{2}\cos\beta) - a_A \\
        a(\omega_{C/A}^{2}\sin\beta - \omega_{C/B}^{2}\sin\gamma) - a_B
    \end{bmatrix}
    $}
\end{equation*}

Una vez resuelto este sistema, es posible determinar las aceleraciones
angulares de ambos brazos y, a partir de ellas, la aceleración absoluta del
punto $C$ \cite{meriam2018}.
